% !TeX root = ../main.tex
% 绪论

\chapter{绪论}

\section{选题背景与研究意义}
    随着软件系统的复杂度不断增加，代码覆盖率已成为评估软件开发质量和效率的关键指标。在嵌入式系统软件开发中，测试需求日益增多且变得更为复杂。例如，在航空航天领域的机载、星载或是弹载软件上的嵌入式软件，就需要具有很高的可靠性以及较强的实时性。进而要求高质量、高效率的嵌入式软件测试来支撑\cite{杨丰玉2021嵌入式软件测试研究综述}。

    代码覆盖率测试就是一种评估嵌入式软件设计质量的有效方法。代码覆盖率定义为在运行特定测试套件时执行/访问程序源代码的程度，可以使用不同的标准来衡量覆盖范围，例如报表覆盖范围、分支机构覆盖范围、条件覆盖范围、MCDC 覆盖范围等\cite{Weiss2024AchievingCS}。通过分析测试过程中源代码的执行程度，能够帮助开发人员发现和解决潜在问题，从而提高软件的可靠性和稳定性。若覆盖率不为100\%，则存在测试用例未覆盖到的代码。这表明测试套件不完整，要么已经实现了测试用例未涵盖的需求，要么已经实现了不必要的功能\cite{Weiss2024AchievingCS}。
    
    嵌入式软件测试与一般的应用软件测试存在显著差异，主要体现在其对硬件的高度依赖性，不能脱离硬件设备独立进行\cite{刘栋2022基于覆盖率的嵌入式软件测试方法研究}。具体到本文的工作中，不同的单板所采用的测试方法在细节上各不相同。

    经过考察，发现部门中代码覆盖率的检视工作存在诸多不足。首先，\textbf{缺乏统一的覆盖率数据解析流程}。不同单板和领域的覆盖率数据解析方式各异，甚至某些领域缺乏自动化解析脚本，导致难以形成联动，全面检视覆盖率。其次，\textbf{覆盖率检视粒度粗}。只能看到单板总体的代码覆盖率数据，无法细分到不同领域或功能模块，测试人员需要检视全量的覆盖率，效率低下且难以定位问题。最后，\textbf{缺乏增量代码覆盖率数据}。无法区分当前版本与历史版本的覆盖率变化，限制了对代码质量的精细化管理和持续改进。

    针对上述问题，本文设计并实现了“多维度覆盖率融合显示系统”。该系统构建了一个通用的自动化覆盖率采集与分析流程，兼容不同单板所使用的交叉编译工具链（LLVM 和 GCC），实现了覆盖率数据处理的统一。系统将代码覆盖率数据划分为多个层次，包括版本级、单板级、领域级及模块级，既能展示版本整体覆盖率，也能细粒度地查看各单板、领域和模块的覆盖率情况。此外，系统支持增量代码覆盖率的提取与展示，能够准确反映不同时间段内代码质量的变化。为提升系统的执行效率和用户体验，避免因接口冲突导致用户排队等待，系统部署在公司内部的 Smart-CI 持续集成平台上，实现了覆盖率分析任务的自动化和分布式执行。

    同时，调研发现，部门现有单板软件测试所依赖的硬件平台采用\textbf{资源受限的异构操作系统架构}，数据面与控制面分别运行于异构 OS 和 Linux OS，其中数据面缺乏独立文件系统，且测试执行结束后会自动断电。这种软件分层与 OS 异构特性，使得覆盖率数据的实时获取与完整保存面临显著挑战。针对这一难题，本文在已有共享内存通信方案的基础上，设计并实现了一套覆盖率数据实时传输机制：在数据面拦截覆盖率统计信息，并通过共享内存传递至控制面，从而实现覆盖率数据的实时、可靠采集，有效避免因断电造成的数据丢失。

    本系统提升了覆盖率数据的可视化能力，使测试人员能够在同一平台上从版本、单板、领域和模块等多个维度全面掌握软件质量状况；统一并优化了覆盖率数据采集与解析流程，减少人工操作，提高了测试工作的自动化程度和整体效率；提供了细粒度的覆盖率检视功能，帮助测试人员快速定位特定功能模块的质量问题；增加了增量覆盖率分析功能，能够准确反映版本迭代过程中的代码质量变化，为项目的持续改进和质量管理提供了有力支持。

\section{国内外研究现状}
    近年来，随着嵌入式系统在工业控制、智能终端和安全关键领域的广泛应用，代码覆盖率分析工具与方法逐渐从通用软件工程扩展到资源受限的嵌入式环境。在此背景下，嵌入式系统通常面临内存与存储受限、实时性要求高、操作系统异构以及缺乏独立文件系统等挑战，这对覆盖率数据的采集、存储和传输提出了更高要求。本节从嵌入式覆盖率分析技术、工程适配方案、资源受限下的覆盖率采集，以及跨核通信机制四个方面，对国内外代表性研究进行综述。

\subsection{覆盖率分析技术}
    目前，主流的代码覆盖率分析技术主要包括静态分析和动态分析两类\cite{张大方2023面向国产嵌入式系统的代码覆盖率实时分析工具}。其一，静态分析在不运行程序的情况下检查源代码逻辑与结构，通常用于处理较小的数据集，依赖于形式化方法和理论基础。适合在开发早期阶段发现死代码、不可达分支及潜在安全漏洞。其二，动态分析则通过跟踪程序运行状态，收集执行路径、变量值及模块交互信息，能够识别仅在实际运行中出现的问题，适用于确保所有代码路径都被执行，并识别代码库中未测试或测试不足的部分。代码覆盖率分析工具和性能分析工具常用于动态分析。

\subsection{LLVM 插桩技术}
    LLVM\cite{llvmoverview} 提供了开源、跨平台的编译框架，其源码级代码覆盖率工具支持在编译过程中插入覆盖率收集指令，从而统计程序各代码段的执行情况。李树芳等人\cite{李树芳2017采用Clang/LLVM的C++源代码覆盖率分析插装方法}提出了一种基于 Clang/LLVM 的C++源代码抽象语法树遍历和代码插桩处理框架，包含了语句覆盖、分支覆盖及 MC/DC 覆盖统计的代码插桩方法。而 Phipps 等人\cite{Phipps2020source}指出，LLVM 插桩技术生成的覆盖率数据量较大，在嵌入式环境中可能造成存储和传输负担，并提出了一种结合硬件辅助采集、数据压缩合并、按需模块采集以及运行时数据过滤的方案，有效平衡了数据完整性与资源开销，使 LLVM 插桩技术能够更加适用于资源受限的嵌入式平台，提升覆盖率分析的可行性和效率。

\subsection{嵌入式环境下的工程适配方案}
    商业化工具在嵌入式领域提供了多维度覆盖率分析支持。Qt Coco\cite{qt2023coco}支持多语言、多平台的代码覆盖率分析，涵盖行覆盖率、分支覆盖率及 MC/DC 覆盖率，并可无缝集成持续集成（CI）和 DevOps 流水线。VectorCAST\cite{vector2020coverage}同样支持多层次覆盖率的采集与分析，覆盖单元测试、集成测试及系统测试等各测试阶段。其白皮书强调，代码覆盖率不仅是提升软件可靠性的关键手段，更是工程实践中系统化测试流程不可或缺的一环。 MPLAB Code Coverage\cite{microchip2020coverage} 以其轻量级设计和与集成开发环境的紧密结合，提供了无需额外硬件支持的高效代码覆盖率采集方案。其高效的数据存储机制显著降低了资源开销，简化了嵌入式系统的测试流程，提升了整体测试效率。
    
    尽管由于封闭的实现细节与商业授权限制了灵活度和可定制性，这些工具仍为本文的系统设计及工程实现提供了先进理念与实践经验，为实现高效、精细的多维度覆盖率管理提供了重要参考。

\subsection{资源受限下的覆盖率采集}
    Weiss 等人\cite{Weiss2024AchievingCS}指出，在嵌入式系统集成测试中，传统方法在目标平台部署测试系统会占用大量资源，并可能引入探测效应，影响覆盖率测量的准确性，并提出了一种通过专用硬件（FPGA）实现非侵入式监控的方案，有效避免了测试系统对运行时的干扰。Saha 等人\cite{Saha2021TraFicA}提出的 TraFic 框架利用系统调试日志和编译器选项来检测代码，降低检测机制的复杂性，实现了低开销的覆盖率采集。在基于ARM Cortex-R 处理器的实验中，TraFic 在略微降低覆盖率目标的情况下，显著降低了性能开销和内存开销。

\subsection{跨核通信技术}
    针对异构嵌入式平台中“Linux 控制面 + 非 Linux 数据面”的场景，跨核通信是实现实时覆盖率采集的关键技术。Coppik 等人\cite{coppik_et_al:OASIcs.NG-RES.2025.2}在 Xilinx Zynq UltraScale+ 异构多处理器系统芯片上，对 IPI、中断+邮箱及共享内存通信进行了实验比较，提出了一种仅依赖共享内存的方案，在小数据量和信号传输场景下可实现亚微秒级延迟。Tsao 等人\cite{tsao10.6688/JISE.2012.28.3.7}评测了 GPP+DSP 双核异构系统中不同 IPC 机制的吞吐、延迟与资源开销，并提出了根据负载和资源约束动态选择通信策略的模型。Hung 等人\cite{hung10.1016/j.sysarc.2010.11.003}提出了名为 MSG 的跨核通信方案，提供阻塞与非阻塞的点对点通信、单边通信和集体操作等接口，并在 IBM CELL 与 ITRI PAC DUO 异构多核平台上进行了实验验证，为嵌入式多核平台提供统一、可移植且高效的通信机制。

\subsection{覆盖率可视化工程实践}
    在实际工程应用中，主流的覆盖率采集与可视化工具通常基于 LCOV\cite{lcov_github}，它能够对 gcov/llvm-cov 生成的执行跟踪数据进行汇总合并，并生成分层次的 HTML 报表，以便开发者进行多维度覆盖率分析。进一步地，Chromium Code Coverage 工具\cite{chromium_coverage} 实现了覆盖率报告与工程流程的深度集成，支持自动化生成细粒度代码覆盖率视图，并将可视化结果集成到代码审查系统 Gerrit 中，提供实时测试反馈和决策支持。这些实践为本文的覆盖率可视化系统设计与工程实现提供了重要参考。

